\subsection{NN-G026: the exact optimum for all constant-ratio closed tables}
\paragraph{Statement.}
The exact cancelling-pair minimum is $4520$ for the closed
rank-three table shape with arbitrary effective cell polynomials
$\delta_1,\delta_2\in\mathbb Z[D]$, arbitrary-degree
$W\in\mathbb Z[D]$, positive rational constant ratios
$G_i/(W\delta_i)$, and arbitrary rational auxiliary functions
$\Lambda_i\in\mathbb Q(D)$. The cells need not be proportional;
there are no degree, monicity, or content restrictions.
Every actual structure constant must still be integral,
effective, coefficientwise dominant over its absolute trace
target, and parity-correct.

The attainer is NN-G025. The new assertion is the exclusion
of every smaller table in this enlarged class. Only the
constant-ratio and closed-shape assumptions remain as
optimization restrictions; an unrestricted dimension-ring
optimum is not asserted.

\paragraph{Normalization without proportionality.}
Write $\delta_i=c_i f_i$, where $f_i\in\mathbb Z[D]$
is primitive effective, and $W=gW_0$ with $W_0$ primitive
effective. All contents $c_i,g$ are positive odd integers.
The reduction of each $f_i$ modulo two is
$\sigma=D(D-1)$, and $W_0\equiv D$.
Gauss's content lemma gives positive integers $a_1$ even,
$a_2$ odd, with $g\mid a_1a_2$, such that
\[
 \begin{aligned}
 G_i&=a_iW_0f_i,&S&=(a_1a_2/g)W_0f_1f_2,\\
 U_i&=\frac{g f_i(a_iW_0-c_iD)}{a_j f_j},&
 V_i&=\frac{g(a_iDW_0f_i-B_i)}{a_j f_j},
 \qquad B_i=\delta_i\Lambda_i.
 \end{aligned}                                                       \tag{G026.1}
\]
Here $B_i$ is the actual integral effective cell coefficient.
Unlike in the proportional case, no integrality of
$c_i\Lambda_i$ is asserted or needed.
Put $d_i=f_i(4)$, $w=W_0(4)$ and $b_i=B_i(4)$.
The masses $d_i,w$ are even and at least four, while
$b_1\ge8$ and $b_2\ge16$.

The exact nonunit mass is
\[
 \begin{split}
 T={}&(a_1a_2/g)wd_1d_2+20w(a_1d_1+a_2d_2)+200gw\\
 &+20\sum_i\{U_i(4)+V_i(4)\}
       +26(c_1d_1+c_2d_2)+14(b_1+b_2).
 \end{split}                                                        \tag{G026.2}
\]
Any improvement would have $T<13464$. The base bounds
$U_i(4)\ge4$, $V_i(4)\ge16$, and $b_1+b_2\ge24$
give the following complete finite ranges:
\[
 \begin{gathered}
 4\le w\le24,\quad4\le d_1\le50,\quad4\le d_2\le78
       \quad\text{all even},\\
 a_1\ge2\text{ even},\quad a_2\ge1\text{ odd},\quad
 a_1+a_2\le33,\quad1\le g\le15\text{ odd},\quad g\mid a_1a_2,\\
 1\le c_i\le125\quad\text{odd}.
 \end{gathered}                                                       \tag{G026.3}
\]
For clarity, these are score bounds, not degree bounds.
Indeed $T\ge472w+1344$ and
$T\ge320(a_1+a_2)+2272$ give the first and scalar-sum
bounds, while $T\ge800g$ gives the content bound for $W$.
Also
\[
 T\ge8d_1d_2+186d_1+106d_2+1936.
\]
Setting $d_2\ge4$ gives $T\ge218d_1+2360$;
setting $d_1\ge4$ gives $T\ge138d_2+2680$.
Finally $T\ge3024+104(c_1+c_2)$ gives an even stronger
bound than the generous individual cap $125$ used here.

\paragraph{An exact scalar relaxation for arbitrary cells and support indices.}
For fixed $a_i,g,w,d_i$, enumerate positive odd $c_i$
within (G026.3) and the necessary bound $c_i\le a_iw/4$.
Impose that
\[
 u_i:=\frac{g d_i(a_iw-4c_i)}{a_jd_j}
              \text{ is an even integer at least four}.             \tag{G026.4}
\]
The coefficient of $c_i$ in (G026.2) is
$d_i(26-80g/(a_jd_j))$. Choose the least admissible content
if that coefficient is nonnegative, and the greatest
otherwise. This gives a lower bound for all contents in
that row. Ignoring their joint bound only weakens the
relaxation, and does not exclude any actual table.

For $b_i$ the necessary condition that $V_i(4)$ be an
even integer at least sixteen is exactly
\[
 \begin{gathered}
 b_i\equiv4a_iwd_i\pmod{M_i},\qquad
 M_i=\frac{2a_jd_j}{\gcd(g,2a_jd_j)},\\
 X_i\le b_i\le
 \left\lfloor\frac{4g a_iwd_i-16a_jd_j}{g}\right\rfloor,
 \qquad (X_1,X_2)=(8,16).
 \end{gathered}                                                       \tag{G026.5}
\]
Its coefficient in (G026.2) is $14-20g/(a_jd_j)$,
so the same endpoint rule minimizes it on this finite
congruence interval. Also impose $G_i(4)\ge48$ and
$S(4)\ge576$. These are all necessary conditions;
no divisibility of $B_i$ by $f_i$ is assumed.

The complete integer-only enumeration
\texttt{verify\_constant\_ratio\_global\_optimum.py}
has $1776$ scalar tuples after the coarse mass bounds,
and exactly $55$ endpoint rows below $T=13464$.
The full JSON certificate lists every row. Thirty-one
have $w=4$. In all remaining twenty-four rows
\[
 (a_1,g)=(2,1),\qquad a_2\in\{1,3\}.
\]
Their mass pairs are exactly
\[
\begin{array}{c|c|l}
 a_2&w&(d_1,d_2)\\\hline
 1&8&(4,12),(6,12),(6,18),(8,16)\\
 1&10&(4,8),(6,8),(6,12),(4,16),(6,16)\\
 1&12&(4,6),(4,8),(4,10),(6,6),(4,12),(6,12)\\
 1&14&(4,8)\\
 1&16&(4,8)\\
 3&8&(4,4),(6,4),(6,6),(10,4)\\
 3&10&(6,4)\\
 3&12&(4,4),(6,4)
\end{array}                                                         \tag{G026.6}
\]
The endpoint contents in the certificate need not equal
those of an actual hypothetical table. The exclusions
below either apply to every content in the row or derive
the actual contents anew from (G026.4).

\paragraph{Three elementary polynomial observations.}
First, coprime $f_1,f_2$ are impossible in (G026.1).
If they were coprime over $\mathbb Q$, polynomiality of
$U_1$ would force $f_2\mid a_1W_0-c_1D$ over $\mathbb Q[D]$.
Primitivity and Gauss's lemma make the quotient integral.
Reducing modulo two would say that $D(D-1)$ divides $D$,
which is false.

Second, when $w=4$, effectivity and parity force $W_0=D$,
and the two cells must be proportional. To see this,
write $f_1=Hu_1$, $f_2=Hu_2$ with primitive integer
polynomials and coprime $u_1,u_2$ over $\mathbb Q$.
The nonzero factors $a_iW_0-c_iD=(a_i-c_i)D$ in
(G026.1) force both $u_1,u_2$ to divide $D$.
A zero factor would make $U_i=0$, which is forbidden.
Since $f_1\equiv f_2\equiv\sigma\ne0$ modulo two,
the nonzero reduction of $H$ gives $u_1\equiv u_2$.
The choices $1$ and $D$ cannot be different, and they
cannot both be $D$ because $u_1,u_2$ are coprime.
Thus both are constant and $f_1=f_2$. NN-G025 excludes
all thirty-one $w=4$ rows below $4520$.

Third, the only effective parity-correct primitive
polynomials of masses four and six are respectively
\[
 f^{(4)}=D(D-3),\qquad f^{(6)}=(D-1)(D-2).
\]
They are coprime. Every such polynomial of mass eight
has the form
\[
 f^{(4)}+4\quad\text{or}\quad
 f^{(4)}+2(q^j+q^{-j})\quad(j\ge1).                   \tag{G026.7}
\]
No polynomial in (G026.7) vanishes at $D=0$ or $D=3$.
For $D=0$, choose $q=-1$, so the added pair has value
$4(-1)^j\ne0$. For $D=3$, choose $q=e^{i\pi/3}$;
the pair contributes $4\cos(j\pi/3)$, which is one
of $4,2,-2,-4$. The constant addition four also
never vanishes. Hence every mass-eight polynomial
is coprime to $f^{(4)}$, uniformly in its unbounded
support location $j$.

\paragraph{Parity modulo four eliminates the nonmultiples of four.}
For the remaining rows $a_1=2$ and $g=1$. Define
\[
 Q=(a_2W_0-c_2D)/2\in\mathbb Z[D],
\]
which is integral because $a_2,c_2$ are odd and
$W_0\equiv D$. The second $U$ identity is
$f_1U_2=f_2Q$. Reducing it modulo two and cancelling
$\sigma\ne0$ gives $Q\equiv D$. Therefore
\[
 a_2W_0\equiv(c_2+2)D\pmod4.                        \tag{G026.8}
\]
At $D=4$ this implies $w\equiv0\pmod4$.
Thus all $w=10$ and $w=14$ rows in (G026.6) are
excluded, regardless of their contents.

\paragraph{The rows with $a_2=1$.}
For $w=8$, positivity of $U_2(4)$ forces $c_2<2$,
so $c_2=1$. Every mass pair listed in this group has
$d_2\ge2d_1$. Then $U_1(4)\ge4$ in (G026.4) forces
$c_1\le2$, hence $c_1=1$ as well.
Equation (G026.8) gives $W_0\equiv3D\pmod4$.
At total mass eight there is exactly one effective
such polynomial:
\[
 W_0=3D-4,
\]
with coefficient list $(2,3)$; any additional outer
pair or increase of a noncentral coefficient modulo
four would already use at least eight further units
of mass.

Write $f_i=Hu_i$ as above. Polynomiality of the two
$U_i$ now forces
\[
 u_1\mid D-2,\qquad u_2\mid5D-8.
\]
Their reductions modulo two agree. Thus either both
are constant, which is the proportional case, or
$u_1=D-2$ and $u_2=5D-8$, up to the already normalized
positive primitive units. In the latter case the
first $U$ identity gives $U_1=D-2$. Its central
Laurent coefficient is zero, below the required
coefficient two of $D$. This excludes all four
$w=8$ rows.

For $w=12$, positivity of $U_2(4)$ gives $c_2<3$,
so again $c_2=1$. Equation (G026.8) requires
$W_0\equiv3D\pmod4$. An effective polynomial of
mass twelve with this residue is uniquely $3D$:
the minimum coefficients are $(2,3)$ of mass eight,
and the remaining four units can only increase the
central coefficient by four. But $3D$ has content
three, contradicting primitivity of $W_0$. All six
rows in this group are excluded.

The only $w=16$ row has mass pair $(4,8)$, whose
primitive cells are coprime by (G026.7). It is
excluded by the first polynomial observation.

\paragraph{The rows with $a_2=3$.}
After the modulo-four test, these have $w=8$ or
$w=12$. Mass pairs $(4,4)$ and $(6,6)$ force identical
primitive cells and are covered by NN-G025; $(6,4)$
is coprime and impossible. The only remaining pair
is $(d_1,d_2)=(10,4)$ at $w=8$.
In this row
\[
 U_1(4)=(10/12)(16-4c_1),\qquad
 U_2(4)=(24-4c_2)/5.
\]
Requiring both to be at least four and the contents
to be positive odd gives $c_1=c_2=1$.
Equation (G026.8) gives $W_0\equiv D\pmod4$.
At mass eight this uniquely gives $W_0=D+4$.
Now $f_2=D(D-3)$ and polynomiality of
\[
 U_1=\frac{f_1(D+8)}{3f_2}
\]
forces $f_2\mid f_1$, because $f_2$ is coprime
to $D+8$. The quotient is integral by Gauss's lemma.
Its value at $D=4$ would be $d_1/d_2=10/4$, which
is not an integer. This last row is impossible.

All fifty-five necessary-relaxation rows have been
excluded below $4520$. The NN-G025 attainer belongs
to the stated class, so the minimum is exactly
$4520$.\hfill$\square$

\paragraph{Geometric dictionary and the smaller remaining gate.}
The parity and modulo-four arguments are constraints
on actual effective graded sign representations:
dimension and signed trace agree modulo two, while
the composing table imposes further integral divisions.
SGA 3, I, 4.7.2--4.7.3 provides the diagonalizable-group
character dictionary, and the NN-G010 realization is
the finite-kernel composition recorded in archive
cote 2, scan 17. The evaluations in (G026.7) restrict
the grading character to finite roots of unity;
they are not Frobenius counts or arithmetic zeta zeros.
No rational auxiliary factor is promoted to an
effective object unless its actual coefficient has
been verified integral and nonnegative.

In geometric terms, all effective graded multiplicity objects
for the two cells may now vary independently, subject to this
closed kernel-table presentation and its constant rational
ratio identities. Among the resulting NN-G010 realizations,
the least total half-excess of dimension over the absolute
signed trace coefficients is exactly $4520$. The trace
specialization is fixed, but no canonical choice of objects,
coherence data, or generators is claimed. This optimum is
presentation-dependent, not a universal categorical invariant.

Arbitrary nonproportional cell polynomials and
unbounded support locations have now been treated
within the constant-ratio closed-table class. A
smaller score must use nonconstant ratios
$G_i/(W\delta_i)$ or a different table shape.
This is an exact realization theorem with explicit
scope, not a canonical-basis theorem or a consequence
for PH0, PH1, rank four, or nonstandard purity.

